\documentclass[11pt]{article}
\usepackage[a4paper,margin=26mm]{geometry}
\usepackage[T1]{fontenc}
\usepackage{lmodern}
\usepackage{amsmath,amssymb,amsthm}
\usepackage{microtype}
\usepackage{array}
\usepackage{xurl}
\usepackage{hyperref}
\hypersetup{colorlinks=true,linkcolor=blue,urlcolor=blue,citecolor=blue,
 pdftitle={Covering minima of the standard terminal simplex and a discrete surface area bound for the covering radius of simplices},
 pdfauthor={Horace Dong}}
\newtheorem{theorem}{Theorem}
\newtheorem{lemma}[theorem]{Lemma}
\newtheorem{proposition}[theorem]{Proposition}
\newtheorem{corollary}[theorem]{Corollary}
\newtheorem{thmprime}{Theorem}
\renewcommand{\thethmprime}{2$'$}
\theoremstyle{definition}
\newtheorem{definition}[theorem]{Definition}
\newtheorem{conjecture}{Conjecture}
\renewcommand{\theconjecture}{\Alph{conjecture}}
\theoremstyle{remark}
\newtheorem{remark}[theorem]{Remark}
\newcommand{\repoversion}{1.25.0}
\setlength{\emergencystretch}{1.5em}
\newcommand{\R}{\mathbb R}
\newcommand{\Z}{\mathbb Z}
\newcommand{\Q}{\mathbb Q}
\newcommand{\one}{\mathbf 1}
\newcommand{\conv}{\operatorname{conv}}
\newcommand{\Vol}{\operatorname{Vol}}
\newcommand{\vol}{\operatorname{vol}}
\newcommand{\Surf}{\operatorname{Surf}}
\newcommand{\fr}[1]{\{#1\}}

\title{Covering minima of the standard terminal simplex\\ and a discrete surface area bound\\ for the covering radius of simplices}
\author{Horace Dong\\ \small Independent researcher\\
\small\texttt{maxwelldhx@gmail.com}\\
\small ORCID: \href{https://orcid.org/0009-0002-7554-7548}{0009-0002-7554-7548}}
\date{22 November 2026}
\begin{document}
\maketitle

\begin{abstract}
The covering radius $\mu(K,\Lambda)$ of a convex body $K$ with respect to a
lattice $\Lambda$ is the least $\mu\ge0$ such that $\mu K+\Lambda$ is the whole
space, and the $i$-th covering minimum $\mu_i(K)$ is the largest covering
radius of a linear projection of $K$ to dimension $i$ under which $\Z^n$ is
mapped onto a lattice. González Merino and Schymura conjectured that the
standard terminal simplex $T_n=\conv\{-\one,e_1,\dots,e_n\}$ satisfies
$\mu_i(T_n)=i/2$ for $1\le i\le n$. Codenotti, Santos and Schymura stated this
as their Conjecture~B, proved it for $i\le3$, and proposed a discrete analogue
of Hadwiger's inequality (their Conjecture~C): every $d$-simplex $S$ with the
origin in its interior and rational vertex directions satisfies
$\mu(S)\le\frac12\Surf(S)/\Vol(S)$, where $\Surf$ is a discrete surface area
obtained by projecting along the vertex directions. Conjecture~C was known for
$d\le2$ and, in general, up to a constant factor. We give proofs of both
conjectures in all dimensions. The proofs are elementary. They average
sawtooth functions over their jump points; for Conjecture~B this is combined
with an induction that splits a rational kernel along the support of a
nonnegative vector, and for Conjecture~C with determinant identities for the
barycentric coordinates of the origin. Both theorems are formally verified in
Lean~4. By published arguments of Codenotti, Santos and Schymura, which are
not formalized here, either theorem implies the inequality $\mu(P)\le d/2$ of
their Conjecture~A for every lattice $d$-polytope $P$ with an interior lattice
point. The equality characterization in Conjecture~A does not follow, and
another paper in this series shows that it fails in dimension four. The proofs
were found and formalized by AI agents working under the author's direction,
as described at the end of the paper; no human mathematician has yet read
them.
\end{abstract}

\section{Introduction}

\subsection{Covering radius and covering minima}
A lattice in $\R^d$ is the image of $\Z^d$ under an invertible linear map.
Following Codenotti, Santos and Schymura~\cite{css}, the covering radius of a
convex body $K\subseteq\R^d$ with respect to a lattice $\Lambda$ is
\[
\mu(K,\Lambda)=\min\{\mu\ge0:\ \mu K+\Lambda=\R^d\},
\]
and $\mu(K)=\mu(K,\Z^d)$. For $K\subseteq\R^n$ and $1\le i\le n$, the $i$-th
covering minimum, introduced by Kannan and Lovász~\cite{kl}, can be written
as~\cite[p.~3]{css}
\[
\mu_i(K,\Lambda)=\max_\pi\ \mu(\pi(K),\pi(\Lambda)),
\]
where $\pi$ runs over the linear maps $\R^n\to\R^i$ for which $\pi(\Lambda)$ is
a lattice in $\R^i$. González Merino and Schymura define $\mu_i$ through affine
subspaces and use the equivalent description by projections onto lattice
planes due to Kannan and Lovász~\cite[eq.~(14)]{gms}; we use the form
of~\cite{css}, which is also the form formalized in Lean. A lattice polytope
is a polytope with vertices in $\Z^d$; it is \emph{non-hollow} if its interior
contains a lattice point. The \emph{standard terminal simplex} is
\[
T_n=S(\one_{n+1})=\conv\{-\one_n,e_1,\dots,e_n\}\subseteq\R^n ,
\qquad \one_n=(1,\dots,1).
\]
We write $\Vol_\Lambda$ for the normalized volume, the Euclidean volume scaled
so that a unimodular simplex of $\Lambda$ has volume one~\cite[p.~3]{css}.
Throughout, $\fr t=t-\lfloor t\rfloor\in[0,1)$ denotes the fractional part of
a real number $t$.

\subsection{The conjectures}
We cite~\cite{gms} by the numbering of arXiv:1605.00443v1, whose second
author is listed there under the name Henze; Codenotti, Santos and Schymura
cite the journal version, apparently with numbers shifted by one: our
Proposition~4.8 and Remark~4.9 are cited in~\cite{css} as Prop.~4.9 and
Rem.~4.10. We could not consult the journal text. The coordinate projection $\R^n\to\R^i$ maps
$T_n$ onto $T_i$ and $\Z^n$ onto $\Z^i$, and $\mu(T_i)=i/2$
\cite[Prop.~4.8]{gms}; hence $\mu_i(T_n)\ge i/2$. González Merino and
Schymura derived this lower bound and conjectured that it is an
identity~\cite[Remark~4.9]{gms}.

\setcounter{conjecture}{1}
\begin{conjecture}[{\cite[Remark~4.9]{gms}; \cite[Conjecture~B]{css}}]\label{conj:B}
For all integers $1\le i\le n$, $\mu_i(T_n)=i/2$.
\end{conjecture}

\noindent (The case $i=0$ is trivial.) Codenotti, Santos and Schymura showed
that Conjecture~B is equivalent to the inequality part of their
Conjecture~A~\cite[Theorem~1.2]{css}:

\begin{quote}
\textbf{Conjecture A of~\cite{css}.} Let $P\subseteq\R^d$ be a non-hollow
lattice $d$-polytope. Then $\mu(P)\le d/2$, with equality if and only if $P$
is obtained by direct sums and/or translations of simplices of the form
$S(\one_l)$.
\end{quote}

\noindent More precisely, \cite[Theorem~1.2]{css} states that, for each $d$,
the inequality $\mu(P)\le\ell/2$ holds for every non-hollow lattice
$\ell$-polytope and every $\ell\le d$ if and only if $\mu_\ell(T_n)=\ell/2$ for
all $\ell\le d\le n$.

For a $d$-simplex $S=\conv\{v_0,\dots,v_d\}$ with the origin in its interior,
$S$ has \emph{rational vertex directions} if each line $\R v_i$ has rational
direction~\cite[Definition~1.5]{css}. Writing $\pi_i\colon\R^d\to\R^{d-1}$ for
a linear projection vanishing at $v_i$, the \emph{discrete surface area} is
\[
\Surf_{\Z^d}(S)=\sum_{i=0}^d\Vol_{\pi_i(\Z^d)}(\pi_i(S))
\quad\text{\cite[Definition~1.5]{css}}.
\]

\setcounter{conjecture}{2}
\begin{conjecture}[{\cite[Conjecture~C]{css}}]\label{conj:C}
Let $S$ be a $d$-simplex with the origin in its interior and with rational
vertex directions. Then
\[
\mu(S)\le\frac12\,\frac{\Surf_{\Z^d}(S)}{\Vol_{\Z^d}(S)} .
\]
\end{conjecture}

\noindent Let $\ell_i$ be the lattice length of the chord $S\cap\R v_i$. Then
$1/\ell_i=\Vol_{\pi_i(\Z^d)}(\pi_i(S))/\Vol_{\Z^d}(S)$, so Conjecture~C reads
$\mu(S)\le\frac12\sum_i1/\ell_i$ \cite[Lemma~4.1, eq.~(6)]{css}. Codenotti,
Santos and Schymura showed that Conjecture~C implies the inequality part of
Conjecture~A~\cite[Corollary~4.3 and its proof]{css}. Corollary~4.3
of~\cite{css} is stated as ``Conjecture~C $\Rightarrow$ Conjecture~A'', but its
proof gives only the inequality; by Theorem~\ref{thm:C} and~\cite[Remark~3]{p21}
the equality clause cannot follow.

\paragraph{What was known.} Conjecture~B holds for $i=n$~\cite[Prop.~4.8]{gms},
for $i=1$ (the lattice width of $T_n$ is $2$; see~\cite[p.~8]{kr}), and for
$i\le3$ and all $n$, since Conjecture~A holds in dimension at most
three~\cite[Theorem~1.3]{css}. Krivokuća~\cite{kr} proved general upper bounds,
for example $\mu_i(T_n)\le\frac i2\bigl(1+\frac{n-i}{n+1}\bigr)$
\cite[Corollary~5.9]{kr} (written there with $d$ for $n$), which equals $7/3$
for $(n,i)=(5,4)$; she lists Conjecture~B as open~\cite[Conjecture~4.2]{kr}.
Conjecture~C holds in dimension two~\cite[Theorem~4.9 and Corollary~4.11]{css};
in every dimension it holds with equality for the simplices
$S(\omega)=\conv\{-\omega_0\one_d,\omega_1e_1,\dots,\omega_de_d\}$
\cite[Theorem~1.4]{css}, it holds for simplices whose primitive vertex
directions sum to zero~\cite[Corollary~5.1]{css}, and it holds in general with the constant $\frac12$
replaced by $1$~\cite[Proposition~4.4]{css} or by
$\frac{2d-1}{2d+2}$~\cite[Remark~4.5]{css}. In dimension one Conjecture~C is
an equality and easy; \cite{css} gives no explicit range of $d$.

\subsection{Results}

\begin{theorem}\label{thm:B}
Let $1\le i\le n$ be integers.
\begin{enumerate}
\item[\textup{(a)}] For every linear surjection $p\colon\R^n\to\R^i$ such that
$p(\Z^n)$ is a lattice, $p(\Z^n)+\frac i2\,p(T_n)=\R^i$.
\item[\textup{(b)}] There is a linear surjection $p\colon\R^n\to\R^i$ such that
$p(\Z^n)$ is a lattice and $i/2$ is the least $r\ge0$ with
$p(\Z^n)+r\,p(T_n)=\R^i$.
\end{enumerate}
Consequently $\mu_i(T_n)=i/2$; that is, Conjecture~B holds.
\end{theorem}

\noindent In the proof of~(b), $p$ is the coordinate projection, as
in~\cite{gms}. Part~(a) says that $\mu(p(T_n),p(\Z^n))\le i/2$ for every
admissible $p$, and (b) that the maximum $i/2$ is attained.

For Conjecture~C we use coordinates. A vector $u\in\Z^d$ is \emph{primitive}
if its coordinates have greatest common divisor~$1$. If $S$ has the origin in
its interior, its vertices are nonzero, and a vertex $v$ with rational
direction can be written uniquely as $v=\omega u$ with $u\in\Z^d$ primitive
and $\omega>0$ real; so the simplices of Conjecture~C are exactly the
simplices of the following theorem. For vectors $w_1,\dots,w_k$ we write
$[w_1;\dots;w_k]$ for the matrix with rows $w_1,\dots,w_k$.

\begin{theorem}\label{thm:C}
Let $d\ge1$, let $u_0,\dots,u_d\in\Z^d$ be primitive, let
$\omega_0,\dots,\omega_d>0$, put $v_i=\omega_iu_i$ and
$S=\conv\{v_0,\dots,v_d\}$, and assume that the origin is an interior point
of~$S$. Put
\[
V=\bigl|\det[(1,v_0);\dots;(1,v_d)]\bigr|,\qquad
D_i=\bigl|\det[u_i;\,v_{s(1)}-v_{s(0)};\dots;v_{s(d-1)}-v_{s(0)}]\bigr|,
\]
where $s(0)<\dots<s(d-1)$ are the indices different from $i$. Then
\[
\Z^d+r\,S=\R^d\qquad\text{with}\qquad r=\frac{D_0+\dots+D_d}{2V}.
\]
\end{theorem}

\noindent The hypothesis implies $V>0$, which the Lean statement also lists.
The number $V$ is $d!\,\vol(S)=\Vol_{\Z^d}(S)$, and
$D_i=\Vol_{\pi_i(\Z^d)}(\pi_i(S))$ for every linear surjection $\pi_i$ with
kernel $\R v_i$ (Lemma~\ref{lem:vol}). Hence $r=\frac12\Surf_{\Z^d}(S)/\Vol_{\Z^d}(S)$,
$\mu(S)\le r$, and Theorem~\ref{thm:C} is Conjecture~C. Primitivity matters
here: replacing $u_i$ by $mu_i$ would multiply $D_i$ by $m$. The proof of
Theorem~\ref{thm:C} does not use primitivity; primitivity is what makes $r$
equal to the bound of Conjecture~C (Lemma~\ref{lem:vol}). Lean also proves
a form of Theorem~\ref{thm:C} in which the volumes are ratios of Lebesgue
measures of the actual projected sets (Theorem~\ref{thm:Cm} in
Section~\ref{sec:C}).

\begin{corollary}[not formalized]\label{cor:A}
For every $d\ge1$, every non-hollow lattice $d$-polytope $P\subseteq\R^d$
satisfies $\mu(P)\le d/2$.
\end{corollary}

\noindent Corollary~\ref{cor:A} follows from Theorem~\ref{thm:B} by the proof
of \cite[Theorem~1.2]{css}, direction (ii)$\Rightarrow$(i), and independently
from Theorem~\ref{thm:C} by the proof of \cite[Corollary~4.3]{css}; both use
the reduction to simplices in~\cite[Corollary~3.2]{css}. We recall the argument
in Section~\ref{sec:cons}. It is not part of the Lean proofs. The bound is
attained by the direct sums of simplices $S(\one_l)$~\cite[eq.~(1) and
Corollary~2.2]{css}. Corollary~\ref{cor:A} is the inequality part of
Conjecture~A only. The ``only if'' part of its equality clause does not
follow, and it fails: the $4$-simplex $\conv\{e_1,e_2,e_3,e_4,-(2,1,1,1)\}$
has covering radius $2$ and normalized volume $6$ and is not obtained from
the simplices $S(\one_l)$ by direct sums and lattice translations~\cite{p21}.
That simplex satisfies Conjecture~C with equality~\cite[Remark~3]{p21}.

\subsection{Method and prior work}
Both proofs rest on an averaging lemma for sawtooth functions
(Lemma~\ref{lem:saw}): for positive integers $q_j$, positive weights $c_j$ and
real $x_j$ there is a real $t$ with
\[
\sum_jc_j\fr{x_j+q_jt}\le\frac12\Bigl(\sum_jc_j-\frac{\sum_jq_jc_j^2}{\sum_jq_jc_j}\Bigr).
\]
It is proved by averaging over the points where some $x_j+q_jt$ is an integer,
using Hermite's identity. The lemma is elementary, averaging arguments of
this kind are classical, and we make no claim of novelty for it.

For Theorem~\ref{thm:B} we write the projected vertices in a basis of the image
lattice. This gives an integer matrix $A$ with $A\one=0$, and the problem
becomes the rounding of a real vector modulo the real kernel of~$A$
(Lemma~\ref{lem:round}). The rounding lemma is proved by induction, splitting
the rational kernel along the support of a nonnegative kernel vector with a
zero entry; the averaging lemma is the base case. The lower bound is the
coordinate projection of~\cite{gms}.

For Theorem~\ref{thm:C}, Cramer's rule gives positive integers $q_i$,
proportional to $\beta_i\omega_i$ where $\beta$ are the barycentric coordinates
of the origin, with $\sum_iq_iu_i=0$, and a cofactor expansion gives
$D_i/V=(1-\beta_i)/\omega_i$. Both facts are close to~\cite{css}: the
primitive positive integer relation among the vertex directions is used
in~\cite[Section~4.3]{css}, of which $(q_i)$ is a positive multiple, and
$(1-\beta_i)/\omega_i=1/\ell_i$ appears in the proof of~\cite[Corollary~4.3]{css},
so that the second identity is the chord form of~\cite[Lemma~4.1]{css}. The
averaging lemma with weights $c_i=1/\omega_i$ then gives exactly
$\frac12\sum_i(1-\beta_i)/\omega_i$. This argument, in chord form, was already
written in the search for problems that preceded the formalization
(Section~\ref{sec:cons}).

As far as we know, Conjecture~B for $4\le i<n$, Conjecture~C for $d\ge3$ and the
bound of Corollary~\ref{cor:A} for $d\ge4$ have not been proved before; the
searches, which cover the literature up to 20 November 2026, are described in
the final section. The closest earlier work on Conjecture~B that we know of
is Krivokuća's paper~\cite{kr}, which narrows the gap between the bounds; the
paper of Codenotti, Freyer and Krivokuća on minimal covering
bodies~\cite{cfk} is related but does not address Conjectures A--C.

\section{An averaging lemma}\label{sec:saw}

\begin{lemma}[Hermite's identity]\label{lem:hermite}
For every real $a$ and every integer $m\ge1$,
$\sum_{k=0}^{m-1}\lfloor a+k/m\rfloor=\lfloor ma\rfloor$, and hence
\[
\sum_{k=0}^{m-1}\fr{a+k/m}=\fr{ma}+\frac{m-1}2 .
\]
\end{lemma}

\noindent The first identity is in Mathlib; the second follows from
$\sum_{k<m}k/m=(m-1)/2$.

\begin{lemma}\label{lem:saw}
Let $J$ be a finite nonempty set, and for $j\in J$ let $q_j\ge1$ be an
integer, $c_j>0$ and $x_j$ real. Put $W=\sum_jq_jc_j$. There is $t\in\R$ with
\[
\sum_{j\in J}c_j\fr{x_j+q_jt}\le\frac12\Bigl(\sum_{j}c_j-\frac1W\sum_{j}q_jc_j^2\Bigr).
\]
In particular, if all $c_j=1$, there is $t$ with
$\sum_j\fr{x_j+q_jt}\le(|J|-1)/2$.
\end{lemma}

\begin{proof}
Let $B(s)=\fr s-\frac12$. For $j\in J$ and $0\le k<q_j$ put
$e_{j,k}=-(x_j+k)/q_j$; these are the points where $x_j+q_jt$ is an integer.
Coincident points with different labels $(j,k)$ are kept separate and counted
with multiplicity. Since $t-e_{j,k}=t+x_j/q_j+k/q_j$, Lemma~\ref{lem:hermite}
gives
\begin{equation}\label{eq:jumps}
B(x_j+q_jt)=\sum_{k=0}^{q_j-1}B(t-e_{j,k})\qquad(t\in\R).
\end{equation}
Let $F(t)=\sum_jc_jB(x_j+q_jt)$ and, for $i,j\in J$,
$E_{ij}=\sum_kB(x_j+q_je_{i,k})$. By~\eqref{eq:jumps},
$E_{ij}=\sum_{k,l}B(e_{i,k}-e_{j,l})$, so
\[
E_{ij}+E_{ji}=\sum_{k,l}\bigl(B(z_{kl})+B(-z_{kl})\bigr),\qquad z_{kl}=e_{i,k}-e_{j,l}.
\]
Now $B(z)+B(-z)$ is $0$ if $z\notin\Z$ and $-1$ if $z\in\Z$; hence
$E_{ij}+E_{ji}\le0$. For $i=j$, $x_i+q_ie_{i,k}=-k$ is an integer, so
$E_{ii}=-q_i/2$. Give the point $e_{i,k}$ the weight $c_i$; the total weight is
$W$. Then
\[
2\sum_{i,k}c_iF(e_{i,k})=2\sum_{i,j}c_ic_jE_{ij}
=\sum_{i,j}c_ic_j(E_{ij}+E_{ji})\le\sum_ic_i^2\cdot2E_{ii}=-\sum_iq_ic_i^2 .
\]
So some point $t=e_{i,k}$ has $F(t)\le-\frac1{2W}\sum_iq_ic_i^2$, and
$\sum_jc_j\fr{x_j+q_jt}=F(t)+\frac12\sum_jc_j$. For $c_j=1$ the bound is
$\frac12(|J|-1)$.
\end{proof}

\section{Proof of Theorem~\ref{thm:B}}\label{sec:B}

\begin{lemma}[kernel rounding]\label{lem:round}
Let $I$ be a finite nonempty set with $m$ elements, let $K\subseteq\Q^I$ be a
subspace of dimension $k$ that contains a vector with all coordinates
positive, and let $H\subseteq\R^I$ be a real subspace containing $K$. Then for
every $x\in\R^I$ there is $h\in H$ with
\[
\sum_{j\in I}\fr{x_j+h_j}\le\frac{m-k}2 .
\]
\end{lemma}

\begin{proof}
By strong induction on $m$. Let $\eta\in K$ have positive coordinates; then
$k\ge1$.

If $k=1$, then $K=\Q\eta$, and clearing denominators gives $w\in K$ with
positive integer coordinates. By Lemma~\ref{lem:saw} with all $c_j=1$ there is
$t$ with $\sum_j\fr{x_j+w_jt}\le(m-1)/2$; take $h=tw$.

Let $k\ge2$ and choose $u\in K$ not proportional to $\eta$. Let
$a=\min_ju_j/\eta_j$ and $\nu=u-a\eta\in K$. Then $\nu\ge0$, some coordinate
of $\nu$ vanishes, and $\nu\ne0$. Let $I_+=\{j:\nu_j>0\}$ and
$I_0=I\setminus I_+$; both are nonempty, so $|I_+|,|I_0|<m$. For a subset
$X\subseteq I$ let $\rho_X$ denote restriction of coordinates to $X$. Put
\[
K_0=\rho_{I_0}(K),\qquad K_+=\rho_{I_+}\bigl(K\cap\ker\rho_{I_0}\bigr),
\]
and define $H_0$ and $H_+$ from $H$ in the same way; then $K_0\subseteq H_0$
and $K_+\subseteq H_+$. The space $K_0$ contains $\rho_{I_0}(\eta)>0$, and
$K_+$ contains $\rho_{I_+}(\nu)>0$. The kernel of $\rho_{I_0}|_K$ is
$K\cap\ker\rho_{I_0}$, on which $\rho_{I_+}$ is injective; by rank--nullity,
$\dim K_0+\dim K_+=k$.

By induction on $I_0$, there is $g\in H$ with
$\sum_{j\in I_0}\fr{x_j+g_j}\le(|I_0|-\dim K_0)/2$. By induction on $I_+$,
applied to $\rho_{I_+}(x+g)$, there is $f\in H$ with $\rho_{I_0}(f)=0$ and
$\sum_{j\in I_+}\fr{x_j+g_j+f_j}\le(|I_+|-\dim K_+)/2$. Then $h=g+f\in H$
agrees with $g$ on $I_0$, and adding the two bounds gives $(m-k)/2$.
\end{proof}

Let $v_0=-\one_n$ and $v_j=e_j$ ($1\le j\le n$) be the vertices of $T_n$, so
that $\sum_jv_j=0$.

\begin{proposition}[upper bound]\label{prop:upper}
Let $1\le i\le n$ and let $p\colon\R^n\to\R^i$ be a linear surjection such
that $p(\Z^n)$ is the lattice generated by a basis $b_1,\dots,b_i$ of $\R^i$.
Then $p(\Z^n)+\frac i2\,p(T_n)=\R^i$.
\end{proposition}

\begin{proof}
Since $v_j\in\Z^n$, we can write $p(v_j)=\sum_kA_{kj}b_k$ with an integer
matrix $A$ of size $i\times(n+1)$. Each row of $A$ sums to zero, that is
$A\one=0$, because $\sum_jv_j=0$. Let $K=\ker_\Q A\subseteq\Q^{n+1}$ and
$H=\ker_\R A$. Then $\one\in K$ and $\dim K\ge n+1-i$.

Let $y\in\R^i$. Since $p$ is onto and $v_1,\dots,v_n=e_1,\dots,e_n$ span $\R^n$, we have $y=p(\sum_jx_jv_j)$
for some $x\in\R^{n+1}$. Lemma~\ref{lem:round} gives $h\in H$ with
$s:=\sum_j\fr{x_j+h_j}\le i/2$. Since $Ah=0$, we have $p(\sum_jh_jv_j)=0$. Put
$z=\sum_j\lfloor x_j+h_j\rfloor v_j\in\Z^n$ and $a_j=\fr{x_j+h_j}\ge0$; then
$y=p(z)+p(\sum_ja_jv_j)$. With $r=i/2$ and
$w_j=\bigl(a_j+(r-s)/(n+1)\bigr)/r$ we have $w_j\ge0$, $\sum_jw_j=1$ and,
since $\sum_jv_j=0$, $r\sum_jw_jv_j=\sum_ja_jv_j$. Hence
$y=p(z)+r\,p(y')$ with $y'=\sum_jw_jv_j\in T_n$.
\end{proof}

\begin{proposition}[lower bound]\label{prop:lower}
Let $1\le i\le n$ and let $p(x)=(x_1,\dots,x_i)$. Then $p(\Z^n)=\Z^i$, and if
$r\ge0$ and $\Z^i+r\,p(T_n)=\R^i$, then $r\ge i/2$.
\end{proposition}

\begin{proof}
The first claim is clear. Let $x\in\R^i$ have coordinates $x_k=k/(i+1)$ and
suppose $x=p(z)+r\,p(y)$ with $z\in\Z^n$ and $y\in T_n$. Write
$y=\sum_{j=0}^nw_jv_j$ with $w_j\ge0$ and $\sum_jw_j=1$, so that
$y_k=w_k-w_0$. Put $a=rw_0$. For $1\le k\le i$ the $k$-th coordinate gives
$a+k/(i+1)=z_k+rw_k$, so $\fr{a+k/(i+1)}=\fr{rw_k}\le rw_k$; also
$\fr a\le a=rw_0$. Summing over $k=0,\dots,i$ and using Lemma~\ref{lem:hermite},
\[
\fr{(i+1)a}+\frac i2=\sum_{k=0}^{i}\fr{a+\tfrac k{i+1}}\le r\sum_{k=0}^iw_k\le r .
\qedhere
\]
\end{proof}

\begin{proof}[Proof of Theorem~\ref{thm:B}]
A lattice in $\R^i$ is the set of integer combinations of a basis, so (a) is
Proposition~\ref{prop:upper}. For (b), the coordinate projection is a
surjection with $p(\Z^n)=\Z^i$; by (a), $i/2$ is a covering scale, and by
Proposition~\ref{prop:lower} no smaller nonnegative scale is. Hence
$\mu(p(T_n),p(\Z^n))\le i/2$ for every admissible $p$, with equality for the
coordinate projection, so $\mu_i(T_n)=i/2$.
\end{proof}

\section{Proof of Theorem~\ref{thm:C}}\label{sec:C}

Let $d$, $u_i$, $\omega_i$, $v_i$, $S$, $V$ and $D_i$ be as in
Theorem~\ref{thm:C}. Since the origin is an interior point of $S$, the
vertices $v_0,\dots,v_d$ are affinely independent, every $y\in\R^d$ is a
combination $\sum_ia_iv_i$, and the barycentric coordinates
$\beta_0,\dots,\beta_d$ of the origin are positive with $\sum_i\beta_i=1$ and
$\sum_i\beta_iv_i=0$. As $d\ge1$, each $\beta_i<1$. Let $M$ be the matrix with
rows $(1,v_0),\dots,(1,v_d)$, so $V=|\det M|>0$.

\begin{lemma}\label{lem:freq}
Let $q_i=|\det[u_j]_{j\ne i}|$, the absolute value of the determinant of the
$d\times d$ matrix with rows $u_j$, $j\ne i$, in increasing order. Then
$q_i=m\beta_i\omega_i$ with $m=V/\prod_j\omega_j$. In particular the $q_i$ are
positive integers and $\sum_iq_iu_i=0$.
\end{lemma}

\begin{proof}
Replacing the row $(1,v_i)$ of $M$ by $(1,0)=\sum_j\beta_j(1,v_j)$ multiplies
the determinant by $\beta_i$. Expanding the new matrix along the replaced row
gives $\pm\det[v_j]_{j\ne i}$. Hence
$|\det[v_j]_{j\ne i}|=\beta_iV$, and
$\det[v_j]_{j\ne i}=\prod_{j\ne i}\omega_j\cdot\det[u_j]_{j\ne i}$. Finally
$\sum_iq_iu_i=m\sum_i\beta_i\omega_iu_i=m\sum_i\beta_iv_i=0$.
\end{proof}

\begin{lemma}\label{lem:proj}
For every $i$, $D_i=V(1-\beta_i)/\omega_i$.
\end{lemma}

\begin{proof}
Let $\zeta=(0,u_i)$. Then $(1,v_i)=\omega_i\zeta+\sum_j\beta_j(1,v_j)$, so by
multilinearity of the determinant in row $i$,
$\det M=\omega_i\det M_{i\leftarrow\zeta}+\beta_i\det M$, where
$M_{i\leftarrow\zeta}$ is $M$ with row $i$ replaced by~$\zeta$. Thus
$|\det M_{i\leftarrow\zeta}|=V(1-\beta_i)/\omega_i$. In
$M_{i\leftarrow\zeta}$, subtract the row $(1,v_{s(0)})$ from the other rows
$(1,v_{s(k)})$ and expand along the first column; up to sign this gives the
determinant of $[u_i;\,v_{s(1)}-v_{s(0)};\dots;v_{s(d-1)}-v_{s(0)}]$.
\end{proof}

\begin{lemma}\label{lem:vol}
$V=\Vol_{\Z^d}(S)$, and $D_i=\Vol_{\pi(\Z^d)}(\pi(S))$ for every linear
surjection $\pi\colon\R^d\to\R^{d-1}$ with $\pi(v_i)=0$.
\end{lemma}

\begin{proof}
The first claim is the usual determinant formula for the volume of a simplex.
Since $u_i$ is primitive, there is $U\in GL_d(\Z)$ with $Uu_i=e_1$. Let
$p_i(x)$ be the vector of the last $d-1$ coordinates of $Ux$. Then $p_i$ is
surjective, $p_i(\Z^d)=\Z^{d-1}$ and $p_i(v_i)=0$. Since
$0=\sum_j\beta_jp_i(v_j)$ with $\beta_j>0$, the point $p_i(v_i)=0$ lies in
$\conv\{p_i(v_j):j\ne i\}$, so $p_i(S)$ is the image of the facet opposite
$v_i$, a simplex with edge vectors $p_i(v_{s(k)}-v_{s(0)})$. Multiplying the
matrix in the definition of $D_i$ on the right by $U^{\mathsf T}$ changes its
determinant only by a sign and turns its top row into $e_1$; expanding along
that row gives
$D_i=|\det[p_i(v_{s(1)}-v_{s(0)});\dots;p_i(v_{s(d-1)}-v_{s(0)})]|
=\Vol_{\Z^{d-1}}(p_i(S))$ (for $d=1$ both sides are~$1$). Any other $\pi$ as
in the statement is $\phi\circ p_i$ for a linear isomorphism $\phi$ of
$\R^{d-1}$, and $\Vol_{\phi(\Lambda)}(\phi(X))=\Vol_\Lambda(X)$.
\end{proof}

\noindent Lean proves Lemma~\ref{lem:vol} for the projections $p_i$ of the
proof, as part of Theorem~\ref{thm:Cm} below; the last sentence of the proof,
the independence of the choice of projection, is not formalized.

\begin{proof}[Proof of Theorem~\ref{thm:C}]
By Lemma~\ref{lem:proj},
\[
r=\frac{\sum_iD_i}{2V}=\frac12\sum_i\frac{1-\beta_i}{\omega_i}>0 .
\]
Let $c_i=1/\omega_i$. By Lemma~\ref{lem:freq}, $W=\sum_iq_ic_i=m\sum_i\beta_i=m$
and $\sum_iq_ic_i^2=m\sum_i\beta_i/\omega_i$, so the bound of
Lemma~\ref{lem:saw} is
$\frac12\bigl(\sum_ic_i-\sum_i\beta_ic_i\bigr)=r$.

Let $y\in\R^d$ and write $y=\sum_ia_iv_i=\sum_ix_iu_i$ with
$x_i=a_i\omega_i$. Lemma~\ref{lem:saw} gives $t$ with
$\sum_ic_i\fr{x_i+q_it}\le r$. As $\sum_iq_iu_i=0$, we have
$y=\sum_i(x_i+q_it)u_i$. Put $z=\sum_i\lfloor x_i+q_it\rfloor u_i\in\Z^d$
and $\alpha_i=c_i\fr{x_i+q_it}\ge0$. Then $y-z=\sum_i\alpha_iv_i$ and
$\sum_i\alpha_i\le r$. With $w_i=\bigl(\alpha_i+(r-\sum_j\alpha_j)\beta_i\bigr)/r$
we get $w_i\ge0$, $\sum_iw_i=1$ and $r\sum_iw_iv_i=y-z$, because
$\sum_i\beta_iv_i=0$. So $y\in z+rS$.
\end{proof}

The following form uses Lebesgue measures of the actual projected sets. Let
$\lambda_k$ be Lebesgue measure on $\R^k$ ($\lambda_0$ is the unit point mass
on $\R^0$) and $\Delta_k=\conv\{0,e_1,\dots,e_k\}$. If $\Lambda=\Z^k$ and
$X\subseteq\R^k$, then $\Vol_\Lambda(X)=\lambda_k(X)/\lambda_k(\Delta_k)$.

\begin{thmprime}\label{thm:Cm}
Under the hypotheses of Theorem~\ref{thm:C} there are linear surjections
$p_0,\dots,p_d\colon\R^d\to\R^{d-1}$ such that, for every $i$, $p_i(v_i)=0$,
$p_i(\Z^d)=\Z^{d-1}$ and $\lambda_{d-1}(p_i(S))/\lambda_{d-1}(\Delta_{d-1})=D_i$,
and
\[
\Z^d+r'S=\R^d\qquad\text{with}\qquad
r'=\frac{\sum_i\lambda_{d-1}(p_i(S))/\lambda_{d-1}(\Delta_{d-1})}{2\,\lambda_d(S)/\lambda_d(\Delta_d)} .
\]
\end{thmprime}

\noindent Since $p_i(\Z^d)=\Z^{d-1}$, the numerator is
$\sum_i\Vol_{p_i(\Z^d)}(p_i(S))$ and the denominator is $2\Vol_{\Z^d}(S)$, so
$r'$ is the right-hand side of Conjecture~C computed with the particular
projections $p_i$.

\section{Consequences and remarks (not formalized)}\label{sec:cons}

The statements in this section are written arguments; none of them is part of
the Lean statements.

\paragraph{Corollary~\ref{cor:A}.} The following argument is that
of~\cite[Section~3.1]{css}; we recall it to make explicit which part of
Theorem~\ref{thm:B} it uses. Let $P\subseteq\R^d$ be a non-hollow lattice
polytope; after a lattice translation the origin is an interior point. By
Carathéodory's theorem, $P$ contains a lattice simplex $S$, spanned by vertices of $P$, of some
dimension $1\le\ell\le d$ (the origin is not a vertex) with the origin in its
relative interior~\cite[Lemma~3.1]{css}.

Suppose first $\ell=d$, so $\mu(P)\le\mu(S)$. Write $S=\conv\{w_0,\dots,w_d\}$.
The barycentric coordinates of the origin are positive and rational, so
$0=\sum_ib_iw_i$ with positive integers $b_i$; let $N=\sum_ib_i\ge d+1$. The
$N$ vertices $-\one,e_1,\dots,e_{N-1}$ of $T_{N-1}$ sum to zero, so there is a
linear map $\pi\colon\R^{N-1}\to\R^d$ sending exactly $b_i$ of them to $w_i$
for each $i$. Then $\pi(T_{N-1})=S$, and $\pi(\Z^{N-1})\subseteq\Z^d$ is a
lattice because $S$ is full-dimensional. Theorem~\ref{thm:B}(a) and the
inclusion of lattices give
$\mu(S,\Z^d)\le\mu(S,\pi(\Z^{N-1}))\le d/2$.

If $\ell<d$, let $L$ be the linear span of $S$ and $\pi'$ the orthogonal
projection with kernel $L$. By~\cite[Lemma~2.1]{css} and $S\subseteq P\cap L$,
$\mu(P)\le\mu(S,\Z^d\cap L)+\mu(\pi'(P),\pi'(\Z^d))$. The first term is at
most $\ell/2$ by the full-dimensional case applied in $L$, and the second is
at most $(d-\ell)/2$ by induction on $d$, since $\pi'(P)$ is a non-hollow
lattice polytope with respect to $\pi'(\Z^d)$ (identifying $\pi'(\Z^d)$ with
$\Z^{d-\ell}$ by a linear isomorphism)~\cite[Corollary~3.2]{css}.

From Theorem~\ref{thm:C} the same bound follows as in the proof
of~\cite[Corollary~4.3]{css}: for a lattice simplex with the origin in its
interior, each $\omega_i$ is a positive integer, so
$\frac12\sum_i(1-\beta_i)/\omega_i\le\frac12\sum_i(1-\beta_i)=d/2$; the
reduction to simplices is the one just described. Conversely, the inequality
of Corollary~\ref{cor:A} in all dimensions implies Conjecture~B by
\cite[Theorem~1.2]{css}, direction (i)$\Rightarrow$(ii). So
Theorem~\ref{thm:C} implies Theorem~\ref{thm:B} through published arguments;
we do not know a deduction in the other direction. The two Lean proofs are
independent of each other except that they share Lemma~\ref{lem:saw}.

\paragraph{Equality.} The equality clause of Conjecture~A is not addressed by
Theorems~\ref{thm:B} and~\ref{thm:C}. Its ``if'' direction was known~\cite[eq.~(1)
and Corollary~2.2]{css}; its ``only if'' direction fails in dimension four by
the example of~\cite{p21}. For that simplex all $\omega_i=1$ and
$\beta=\frac16(2,1,1,1,1)$, so $\frac12\sum_i(1-\beta_i)=2=\mu$: it is an
equality case of Conjecture~C that is not unimodularly equivalent to any
$S(\omega)$~\cite[Remark~5]{p21} (other equality cases are noted
after~\cite[Corollary~5.1]{css}). Corollary~\ref{cor:A}
should therefore not be read as a proof of Conjecture~A, only of its
inequality. Conjecture~D of~\cite{css}, on polytopes with $k$ interior lattice
points, is not addressed either, nor is the extension of Conjecture~C to
simplices with the origin on the boundary~\cite[Section~4.4]{css}.

\paragraph{Other written arguments.} The search for problems in which these
conjectures were selected also produced written arguments. For Conjecture~B it
gave a proof by averaging over the common zeros on a torus of systems of
integer linear forms, with determinant weights; the proof of
Theorem~\ref{thm:B} given here, by kernel splitting, is different and does not
use it. For Conjecture~C a note of that search already contained the argument
of Section~\ref{sec:C} in chord form: the weighted averaging over the jump
points (with a gcd computation in place of Hermite's identity), the positive
integer relation among the vertex directions, and the identity
$1/\ell_i=(1-\beta_i)/\omega_i$. A review note and a second note of the
search gave an integral version of the averaging and, in the second note, with
all weights equal to~$1$, the bound $\mu(S)\le d/2$ for non-hollow lattice
simplices. The determinant form of Lemmas~\ref{lem:freq}--\ref{lem:vol},
Theorem~\ref{thm:Cm} and the Lean proofs come from the later session. These
notes have not been reviewed for this paper and are not claims of this paper.

{\sloppy
\section{Formal verification}\label{sec:lean}

Theorems~\ref{thm:B}, \ref{thm:C} and~\ref{thm:Cm} are formally verified in
Lean~4 (version 4.34.1) with Mathlib~\cite{mathlib} (commit
\texttt{d13f23b7}). The deduction of $\mu_i(T_n)=i/2$ from (a) and (b) and of
$\mu(S)\le r$ from the covering statement is immediate from the definition of
$\mu$ as a minimum. The identification of $V$ and $D_i$ with the volumes
of~\cite{css} is formalized in Theorem~\ref{thm:Cm} and
\texttt{MeasureModel.normalized\_\allowbreak measure}, except for the
independence of the choice of projection (Lemma~\ref{lem:vol}).
Corollary~\ref{cor:A} and Section~\ref{sec:cons} are not formalized. In
version~\repoversion{} of the repository~\cite{repo} the files are in the
directories \texttt{lean/Research/TaskS/B/} ($11$ files),
\texttt{lean/Research/TaskS/C/} ($13$ files) and
\texttt{lean/Research/TaskS/Shared/} (\texttt{Sawtooth.lean}, used by both);
records of the freezing and of the checks are in the folder
\texttt{certificates/css-covering-minima}.

\paragraph{The statement for Conjecture~B.}
The statement file \texttt{B/Challenge.lean} contains only definitions and
statements, no proofs. In ASCII transliteration (the docstrings and the
\texttt{import}, \texttt{open}, \texttt{namespace} and \texttt{end} lines
omitted; the file itself uses Unicode, e.g.\ $\mathbb R$, $\to$, $\forall$,
$\exists$, $\in$, $\le$, $\wedge$, $\sum$, $\bullet$):
{\footnotesize
\begin{verbatim}
abbrev Space (n : Nat) := Fin n -> Real
abbrev Lattice (n : Nat) := Fin n -> Int
def embed {n : Nat} (z : Lattice n) : Space n := fun j => z j
noncomputable def vertex (n : Nat) : Fin (n + 1) -> Space n :=
  Fin.cases (fun _ => -1) (fun j => Pi.single j 1)
noncomputable def terminal (n : Nat) : Set (Space n) :=
  convexHull Real (Set.range (vertex n))
def LatticeProjection {n i : Nat} (p : Space n ->l[Real] Space i) : Prop :=
  Function.Surjective p /\ exists b : Module.Basis (Fin i) Real (Space i),
    Set.range (fun z : Lattice n => p (embed z)) =
      Set.range (fun w : Lattice i => sum j, (w j : Real) * b j)
def Covers {n i : Nat} (p : Space n ->l[Real] Space i) (r : Real) : Prop :=
  forall x : Space i, exists z : Lattice n, exists y in terminal n,
    x = p (embed z) + r * p y
def ConjectureB : Prop :=
  forall (n i : Nat), 0 < i -> i <= n ->
    (forall p : Space n ->l[Real] Space i, LatticeProjection p -> Covers p ((i : Real) / 2)) /\
    exists p : Space n ->l[Real] Space i, LatticeProjection p /\
      IsLeast {r : Real | 0 <= r /\ Covers p r} ((i : Real) / 2)
\end{verbatim}}
\noindent Here \texttt{->l[Real]} is the type of real-linear maps (written
$\to_{\ell}[\mathbb R]$ in the file), \texttt{*} between a scalar and a vector
stands for the scalar multiple $\bullet$, \texttt{sum j} is a sum over
$\{0,\dots,i-1\}$, and \texttt{IsLeast A c} says that $c\in A$ and $c\le r$
for all $r\in A$. \texttt{vertex n} lists $-\one_n,e_1,\dots,e_n$ and
\texttt{terminal n} is $T_n$. \texttt{LatticeProjection p} says that $p$ is
onto and that $p(\Z^n)$ is the set of integer combinations of some real basis
of $\R^i$, that is, a lattice; every lattice, with any basis, is allowed, and
dense images such as that of $(x,y)\mapsto x+\sqrt2\,y$ are excluded
(surjectivity follows from the second clause). \texttt{Covers p r} says that
$p(\Z^n)+r\,p(T_n)=\R^i$. \texttt{ConjectureB} is Theorem~\ref{thm:B}(a) and~(b)
for all $1\le i\le n$. The final reviewers checked in particular that the
universal bound ranges over all admissible $p$ while the witness appears only
in the attainment clause; the stronger reading that every admissible
projection has covering radius $i/2$ would be false, since
$(x_1,x_2)\mapsto x_1+x_2$ maps $T_2$ onto $[-2,1]$, of covering radius
$1/3$. One reviewer restated \texttt{ConjectureB} fully unfolded, using only
Mathlib's standard notions, and derived it from the checked theorem.

\paragraph{The statement for Conjecture~C.}
In the same transliteration, \texttt{C/Challenge.lean} reads ($\omega$ is
written \texttt{omega}):
{\footnotesize
\begin{verbatim}
abbrev Space (d : Nat) := Fin d -> Real
abbrev Lattice (d : Nat) := Fin d -> Int
def embed {d : Nat} (u : Lattice d) : Space d := fun k => (u k : Real)
def Primitive {d : Nat} (u : Lattice d) : Prop :=
  exists a : Lattice d, sum k, a k * u k = 1
def vertices {d : Nat} (u : Fin (d+1) -> Lattice d) (omega : Fin (d+1) -> Real)
    (j : Fin (d+1)) : Space d := omega j * embed (u j)
noncomputable def simplex {d : Nat} (v : Fin (d+1) -> Space d) : Set (Space d) :=
  convexHull Real (range v)
noncomputable def volume {d : Nat} (v : Fin (d+1) -> Space d) : Real :=
  |(Matrix.of (fun i => Fin.cons (1 : Real) (v i))).det|
noncomputable def projectedVolume {n : Nat} (v : Fin (n+2) -> Space (n+1))
    (u : Fin (n+2) -> Lattice (n+1)) (i : Fin (n+2)) : Real :=
  |(Matrix.of (Fin.cons (embed (u i)) (fun k : Fin n =>
    v (i.succAbove k.succ) - v (i.succAbove 0)))).det|
def Covers {d : Nat} (K : Set (Space d)) (r : Real) : Prop :=
  forall x : Space d, exists z : Lattice d, exists y in K, x = embed z + r * y
def conjectureC : Prop :=
  forall (n : Nat) (u : Fin (n+2) -> Lattice (n+1)) (omega : Fin (n+2) -> Real),
    (forall i, Primitive (u i)) -> (forall i, 0 < omega i) ->
    let v := vertices u omega
    0 < volume v -> (0 : Space (n+1)) in interior (simplex v) ->
    Covers (simplex v) ((sum i, projectedVolume v u i) / (2 * volume v))
\end{verbatim}}
\noindent Here \texttt{range} is \texttt{Set.range}, \texttt{a k * u k} is a
product of integers, \texttt{omega j * embed (u j)} and \texttt{r * y} are
scalar multiples, \texttt{|...|} is the absolute value, and
\texttt{interior} is the topological interior in $\R^d$. \texttt{Primitive u}
says that the coordinates of $u$ generate the unit ideal of $\Z$, that is,
$u$ is primitive. \texttt{Matrix.of (fun i => Fin.cons 1 (v i))} is the matrix
with rows $(1,v_i)$, so \texttt{volume v} is $V$; this \texttt{volume} is the
file's own definition, not Mathlib's measure. In \texttt{projectedVolume},
\texttt{i.succAbove} enumerates the indices other than $i$ in increasing
order, and the matrix has top row $u_i$ followed by the rows
$v_{s(k)}-v_{s(0)}$, so \texttt{projectedVolume v u i} is $D_i$.
\texttt{conjectureC} is Theorem~\ref{thm:C} for $d=n+1$, so every $d\ge1$ is
covered. The final reviewers noted the following points. (i)~The
parametrization $v_i=\omega_iu_i$ with primitive $u_i$ and real $\omega_i>0$
gives exactly the simplices of Definition~1.5 of~\cite{css}, including
irrational vertex lengths. (ii)~Primitivity is needed: with merely integral
$u_i$ the right-hand side could be larger. (iii)~The hypothesis
\texttt{0 < volume v} follows from the interior hypothesis and restricts
nothing. (iv)~For $d=1$ each $D_i$ is $1$, the normalized volume of a point,
and the bound is an equality. (v)~\texttt{Covers (simplex v) r} implies
$\mu(S)\le r$; one reviewer derived the form with the infimum of covering
scales in Lean, and another proved by \texttt{Iff.rfl} that
\texttt{conjectureC} unfolds to the statement written out by hand.

\paragraph{The measured form.} Theorem~\ref{thm:Cm} is the theorem
\texttt{TaskS.C.MeasureModel.measured\_\allowbreak conjectureC} in
\texttt{ProjectionModel.lean}, stated for the hypotheses of
\texttt{conjectureC}. Its conclusion asserts the existence of linear maps
\texttt{p i} that are surjective, vanish at $v_i$ and map the image of $\Z^d$
onto the image of $\Z^{d-1}$, with
\texttt{volume (p i \textquotesingle\textquotesingle{} S) / volume (unit n) = ENNReal.ofReal (projectedVolume v u i)},
and the covering of $\R^d$ at the scale in which the sum of the real parts
(\texttt{toReal}) of these quotients is divided by twice that of
\texttt{volume S / volume (unit (n+1))}. Here \texttt{volume} is Lebesgue
measure (with values in $[0,\infty]$) and \texttt{unit k} is $\Delta_k$.

\paragraph{Review and freezing.}
Both statement files were written by Codex agents. The reading of the sources
(\cite{css}, and~\cite{gms} for Conjecture~B) was checked by a Codex agent that
had written neither statement and had not taken part in the earlier search for
problems; its report is dated 20 November 2026, 22:29:40~UTC. The same agent
reviewed both statement files and approved them for freezing at
22:30:07~UTC, after checking that each compiled without diagnostics. Both were
frozen by recording their SHA-256 hashes (\texttt{458363b3\ldots} for
Conjecture~B, \texttt{607e2915\ldots} for Conjecture~C) at 22:30:34~UTC. An
earlier draft of the statement for Conjecture~B had failed to compile at
22:28:45~UTC (the identifier \texttt{Basis} needed the prefix
\texttt{Module.}); the frozen file compiled at 22:29:23~UTC, and the statement
for Conjecture~C at 22:28:56~UTC. According to the development logs, no proof
file was compiled before 22:30:49~UTC, after the freeze. The statement files
were not changed afterwards.

\paragraph{The proofs.}
\texttt{Proof.ConjectureB} combines \texttt{Upper.upper} (part~(a)) with
\texttt{exists\_\allowbreak projection\_\allowbreak lower} (part~(b)).
\texttt{Proof.conjectureC} applies \texttt{geometry} to
\texttt{TaskS.Shared.sawtooth\_\allowbreak weighted} with weights
$1/\omega_i$. Table~\ref{tab:lean} lists the Lean counterparts of the steps of
the written proofs; names are relative to the namespaces \texttt{TaskS.B},
\texttt{TaskS.C} and \texttt{TaskS.Shared}, and ``private'' marks private
declarations. The $12$ files used for Theorem~\ref{thm:B} have $636$ lines
in all, and the $13$ files of \texttt{TaskS/C} have $754$ lines.

\begin{table}[ht]
\centering\footnotesize
\renewcommand{\arraystretch}{1.05}
\begin{tabular}{@{}>{\raggedright\arraybackslash}p{0.36\textwidth}>{\raggedright\arraybackslash}p{0.60\textwidth}@{}}
\hline
Step of the written proof & Lean (file) \\
\hline
Lemma~\ref{lem:hermite} & \texttt{grid\_sum} (\texttt{Shared/Sawtooth}), from
Mathlib's \texttt{Int.sum\_floor\_add\_div} \\
Lemma~\ref{lem:saw} & \texttt{sawtooth\_weighted}, \texttt{sawtooth\_unweighted};
private \texttt{event}, \texttt{grid\_deviation}, \texttt{dev}, \texttt{dev\_diag},
\texttt{dev\_pair} (\texttt{Shared/Sawtooth}) \\
\hline
Lemma~\ref{lem:round}, case $k=1$ & \texttt{Rounding.integer\_direction}
(\texttt{RationalLine}) \\
Lemma~\ref{lem:round}, splitting and rank--nullity &
\texttt{Rounding.positive\_split}, \texttt{Rounding.part}, \texttt{Rounding.supported},
\texttt{Rounding.part\_rank} (\texttt{Split}) \\
Lemma~\ref{lem:round} & \texttt{Rounding.round}, with the one-dimensional input
\texttt{Rounding.LineBound} as a hypothesis (\texttt{KernelRounding}); it is
discharged by \texttt{sawtooth\_unweighted} in \texttt{Upper.round\_matrix}
(\texttt{Balanced}) \\
Proposition~\ref{prop:upper}: the matrix $A$, $A\one=0$, $Ah=0$, the lattice point $z$ &
\texttt{Coordinates.decompose}, \texttt{Coordinates.intcoords},
\texttt{Coordinates.balanced}, \texttt{Coordinates.annihilates},
\texttt{Coordinates.lift\_int} (\texttt{BasisCoordinates}) \\
Proposition~\ref{prop:upper}: $\dim K\ge n+1-i$; padding &
\texttt{Upper.round\_matrix} (\texttt{Balanced}); \texttt{Hull.sum\_vertex},
\texttt{Hull.scaled} (\texttt{Hull}) \\
Proposition~\ref{prop:upper} & \texttt{Upper.upper} (\texttt{Upper}) \\
Proposition~\ref{prop:lower} & \texttt{exists\_projection\_lower}; private
\texttt{coordProjection}, \texttt{coord\_admissible}, \texttt{terminal\_weights},
\texttt{fract\_le\_id} (\texttt{Lower}) \\
\hline
Barycentric coordinates and spanning & \texttt{coordinates}, \texttt{origin\_row}
(\texttt{Coordinates}) \\
Lemma~\ref{lem:freq} & \texttt{frequency}, \texttt{frequency\_ratio}; private
\texttt{replaced\_minor} (\texttt{Frequencies}) \\
Lemma~\ref{lem:proj} & \texttt{projected\_eq}; private \texttt{facet\_det},
\texttt{updated\_facet} (\texttt{Projected}) \\
Rounding and padding in the proof of Theorem~\ref{thm:C} &
\texttt{rounded\_cover}, \texttt{scaled\_combination} (\texttt{Lift}) \\
Proof of Theorem~\ref{thm:C} & \texttt{geometry} (\texttt{Geometry}),
\texttt{Proof.conjectureC} (\texttt{Proof}) \\
Lemma~\ref{lem:vol} for $p_i$: $U$, $p_i$, $p_i(\Z^d)=\Z^{d-1}$ &
\texttt{ProjectionLattice.exists\_adapted} (\texttt{ProjectionLattice});
\texttt{MeasureModel.project}, private \texttt{projection\_surj},
\texttt{projection\_lattice} (\texttt{ProjectionModel}) \\
Lemma~\ref{lem:vol} for $p_i$: image of $S$, expansion, measures &
\texttt{image\_simplex\_eq\_image\_opposite} (\texttt{ProjectionBody});
\texttt{ProjectionLattice.edge\_determinant},
\texttt{ProjectionLattice.projected\_volume};
\texttt{MeasureModel.volume\_edges}, \texttt{MeasureModel.measure\_simplex},
\texttt{MeasureModel.normalized\_measure} (\texttt{SimplexMeasure}) \\
Theorem~\ref{thm:Cm} & \texttt{MeasureModel.projection\_measure},
\texttt{MeasureModel.measured\_conjectureC} (\texttt{ProjectionModel}) \\
\hline
\end{tabular}
\caption{Lean counterparts of the steps of the proofs (files in
\texttt{TaskS/B}, \texttt{TaskS/C} and \texttt{TaskS/Shared}).}
\label{tab:lean}
\end{table}

\noindent The test file \texttt{CheckB.lean} contains
\texttt{theorem check\_\allowbreak ConjectureB : Challenge.ConjectureB :=
Proof.ConjectureB}, and \texttt{CheckC.lean} and
\texttt{MeasureAudit/CheckC.lean} contain the corresponding line
\texttt{check\_\allowbreak conjectureC} for \texttt{Challenge.conjectureC}.
No file other than the statement files defines anything in the namespaces
\texttt{TaskS.B.Challenge} and \texttt{TaskS.C.Challenge}. The axioms of
\texttt{Proof.ConjectureB}, \texttt{check\_\allowbreak ConjectureB},
\texttt{Proof.conjectureC}, \texttt{check\_\allowbreak conjectureC},
\texttt{measured\_\allowbreak conjectureC} and of the other declarations
audited in the checks are exactly \texttt{propext}, \texttt{Classical.choice}
and \texttt{Quot.sound}. No file uses \texttt{sorry},
\texttt{native\_\allowbreak decide}, \texttt{bv\_\allowbreak decide},
\texttt{decide} or an \texttt{axiom} declaration.

\paragraph{Checks.}
Each formalization was compiled module by module, with the pinned toolchain
and with warnings treated as errors, in two new empty directories on the
author's Mac, and every module was replayed with \texttt{leanchecker}; each
directory also ran one build of the whole package. Theorem~\ref{thm:B}
($12$ modules: the $11$ files of \texttt{B} and \texttt{Sawtooth.lean}) was
checked on 20 November 2026 between 22:34:18 and 22:39:44~UTC ($50$
commands), Theorem~\ref{thm:C} ($9$ modules) between 22:33:22 and
22:37:34~UTC ($38$ commands), and Theorem~\ref{thm:Cm} ($14$ modules, which
include the $9$ unchanged modules of Theorem~\ref{thm:C}) between 22:38:51
and 22:44:28~UTC ($58$ commands). All $146$ commands ended with exit
code~$0$, the \texttt{leanchecker} runs printed nothing, and in each check
the two directories produced identical \texttt{.olean} files and axiom lists.
The SHA-256 hashes of the three receipts begin with \texttt{a3307acd},
\texttt{df9cf4b1} and \texttt{4f37b409}. Each of the four final reviewers
(see the last section) compiled the files again, together with an audit file
of its own, in a new directory on the same Mac with plain \texttt{lean}
invocations, and replayed every module with \texttt{leanchecker}; all runs
ended with exit code~$0$ and no warnings (one reviewer's audit file first
failed on an ambiguous name in its own text and was corrected). Mathlib and its
dependencies were taken from a shared prebuilt cache at the pinned commits and
were not rechecked; the \texttt{leanchecker} replays cover the project
modules. In the repository the files have module names under
\texttt{Research.TaskS} and differ from the checked files only in their
\texttt{import} lines (the two statement files, \texttt{Split.lean},
\texttt{RationalLine.lean} and \texttt{Sawtooth.lean} are unchanged).
The repository versions were made by a script that changed only the local \texttt{import} lines, and on 20 November 2026 (23:11--23:18~UTC) the $25$ modules were compiled and replayed with \texttt{leanchecker} again, module by module, in two new empty directories on the author's Mac, one after the other; there were no warnings, and the two directories produced identical \texttt{.olean} files and the same axiom lists. In the whole-project rebuild of version~\repoversion{} of the repository (on the author's Windows computer, in two passes on 20 and 21 November 2026), these files were compiled and replayed with \texttt{leanchecker} again in pass~2 (21 November 2026), which built the files not included in pass~1 on top of its results; no module needed a second attempt (details in the repository). The formalization has not been checked with
comparator or with an independent kernel.

\par}

\section*{How this paper was produced}
This work was carried out from 20 to 22 November 2026 (UTC) in a workflow directed by
the author, using OpenAI Codex with the model GPT-6.2-Astra (reasoning level
ultra) and Anthropic Claude Code with the model Claude Opus~5.5. The author
asked for a search for hard and substantial open problems, and Claude Code
passed this request to Codex. In a separate Codex session, which ended at
18:53~UTC, Codex screened $78$ candidate problems and prepared detailed notes
on eight; Conjectures~B and~C were among them. According to the file times,
the written arguments of this search described in Section~\ref{sec:cons} were
completed by 18:41~UTC (Conjecture~B, torus route), 18:46~UTC (Conjecture~C,
the weighted averaging used in this paper, in chord form; a review note of
18:46~UTC added the integral version) and 18:50~UTC (the second note with the
integral version), and were read by other Codex agents of the search; the counterexample of~\cite{p21}
was found in the same search. At about 22:27~UTC the author asked for
complete Lean proofs of Conjectures~B and~C. In another separate Codex session,
Codex agents (same model) wrote the two statement files, checked the reading
of the sources, searched again for prior work, wrote the Lean proofs and the
written proofs, and ran the checks; the review and freezing of the statements
are described in Section~\ref{sec:lean}. According to the development logs
and the file times, all proof files were created after the freeze; the
modules of Theorems~\ref{thm:B} and~\ref{thm:C} were compiled during
development between 22:30:49 and 22:34:08~UTC and the additional modules of
Theorem~\ref{thm:Cm} between 22:35 and 22:39~UTC. For Conjecture~B
the Lean route differs from the earlier written proof; for Conjecture~C it
follows the earlier note, with the volumes identified through determinants.
Codex's written proofs of the Lean routes, from which Sections~\ref{sec:saw}--\ref{sec:C}
were written, date from the same period (file times 22:30--22:40~UTC).
Because of the significance of the results, each of the two theorems was given
two independent final reviews by Claude Code agents (model Claude Opus~5.5),
four reviewers in all, none of which had written any of the files; each
second reviewer read the first report only after writing its own analysis.
The reviews took place between 22:43 and 23:16~UTC. Each reviewer read the
sources itself, compared the statement file with them, read the written
proof, the proof files and the raw logs of the checks, rebuilt the
formalization as described above; the reviewers between them tested the
statements numerically: the averaging lemma in exact rational arithmetic on
random instances (two reviews); for Conjecture~B, covering radii of sampled
projections with $i\le3$ in one review and of random and adversarially chosen
projections with $n\le5$ and $i\le4$, more than $11{,}000$, in the other,
including the cases $(n,i)=(5,3)$ and $(5,4)$, which earlier results we know
of do not cover; for Conjecture~C, the expressions for the bound (the
determinant formula and the chord form in both reviews, and Definition~1.5
of~\cite{css} with orthogonal projections in one) and lower estimates of
$\mu(S)$ for random simplices in dimensions $2$ to $4$. No
violation was found. All four decisions were APPROVE\_FORMAL, with no
blocking or major remarks; their minor remarks on the written account
($4$ and $5$ for Conjecture~B, $5$ and $7$ for Conjecture~C) are taken into
account here, among them the statement of the consequence for Conjecture~A
with its limits, the numbering of~\cite{gms}, the role of primitivity, and
the status of Lemma~\ref{lem:saw}. Claude Code (Claude Opus~5.5) checked the
sources, checked Lemmas~\ref{lem:saw}, \ref{lem:freq} and~\ref{lem:proj}
numerically in Python, and wrote this text from Codex's manuscripts and the
Lean files. An independent Claude Code agent (model Claude Opus~5.5) that had written none of the files reviewed this paper, and another reviewed the release; their corrections were applied. Throughout, ``we'' refers to this workflow; the
reading and searching described in the paper were done by the agents. Codex
agents searched for prior work in the search for problems and again before
the statements were frozen, and the final reviewers searched again; up to
20 November 2026 these searches found no earlier proof of Conjecture~B for
$4\le i<n$, of Conjecture~C for $d\ge3$, or of the bound of
Corollary~\ref{cor:A} for $d\ge4$ (details in the repository). One citation
index refused the queries for works citing~\cite{css} (rate limit), so a
further search of works citing~\cite{css}, \cite{gms} and~\cite{kr} was
planned for the release. Claude Code made this search on 21 November 2026 (17:54--18:21~UTC), together with searches of the arXiv, Crossref, DataCite, Zenodo, MathDB, GitHub and repositories of AI-assisted mathematics: the same citation index again refused to list the works citing~\cite{css}, Semantic Scholar listed the same ten citing works as before, the works citing~\cite{gms} and~\cite{kr} were unchanged since 20 November, and no new related work was found. Claude Code searched again on 22 November 2026 (05:39--05:51~UTC), when the OpenAlex lists of citing works refused on 21 November were retrieved, and found nothing new. This reports the coverage of the
searches, not a guarantee that they were complete. The author is not a
professional mathematician. The author has read this paper and, using a side-by-side table, compared its Lean statements with the main theorems: Theorems~\ref{thm:B}, \ref{thm:C} and~\ref{thm:Cm}. No human mathematician has read the
proofs or reviewed this paper. The correctness of Theorems~\ref{thm:B},
\ref{thm:C} and~\ref{thm:Cm} rests on the Lean proofs and on the agreement of
the Lean statements with the conjectures of~\cite{gms,css}; the written proofs
are given for the reader. Corollary~\ref{cor:A}, the consequence for
Conjecture~A, is not formalized and rests in addition on the published
arguments of~\cite{css}. AI systems are not authors of this paper; the author
takes full responsibility for its content.

\medskip
\noindent Thanks are due to Bernardo González Merino and Matthias Schymura,
and to Giulia Codenotti, Francisco Santos and Matthias Schymura, for the
conjectures.

{\small
\begin{thebibliography}{9}
\setlength{\itemsep}{1pt}
\bibitem{cfk} G.~Codenotti, A.~Freyer and K.~Krivokuća, Minimal covering
bodies and Brunn--Minkowski type inequalities for the covering radius,
arXiv:2606.14442, 2026.
\bibitem{css} G.~Codenotti, F.~Santos and M.~Schymura, The covering radius and
a discrete surface area for non-hollow simplices, \emph{Discrete Comput.
Geom.} 67 (2022), no.~1, 65--111,
\url{https://doi.org/10.1007/s00454-021-00330-3}; arXiv:1903.02866v2 (2021).
\bibitem{p21} H.~Dong, A non-hollow lattice $4$-simplex of covering radius $2$
outside the equality list of a conjecture of Codenotti, Santos and Schymura,
paper in~\cite{repo}, version~\repoversion{} (2026).
\bibitem{repo} H.~Dong, \emph{ai4math-results}, repository,
\url{https://github.com/Horace-Maxwell/ai4math-results}, version~\repoversion{}
(2026); archived at Zenodo under the concept DOI
\url{https://doi.org/10.5281/zenodo.22970254}.
\bibitem{gms} B.~González Merino and M.~Schymura, On densities of lattice
arrangements intersecting every $i$-dimensional affine subspace,
\emph{Discrete Comput. Geom.} 58 (2017), no.~3, 663--685,
\url{https://doi.org/10.1007/s00454-017-9911-x}; arXiv:1605.00443v1 (2016),
with author line ``B.~González Merino and M.~Henze''.
\bibitem{kl} R.~Kannan and L.~Lovász, Covering minima and lattice-point-free
convex bodies, \emph{Ann. of Math. (2)} 128 (1988), no.~3, 577--602,
\url{https://doi.org/10.2307/1971436}.
\bibitem{kr} K.~Krivokuća, Upper bounds on covering minima of convex bodies,
arXiv:2601.15173v1, 2026.
\bibitem{mathlib}
The mathlib Community, The Lean mathematical library, in \emph{Proc.\ 9th ACM
SIGPLAN Int.\ Conf.\ Certified Programs and Proofs} (CPP '20), ACM, 2020,
\url{https://doi.org/10.1145/3372885.3373824}.
\end{thebibliography}}
\end{document}
